\documentclass[12pt, twoside]{article}
\input{../preamble}

\fancyhead[LO]{La Scuola d'Italia / Dr.\ Huson / IB Math 12 \\* 9 September 2026}
\fancyhead[RO]{\fbox{\begin{minipage}{6cm}\footnotesize First \& Last name:\\[0.5cm]\end{minipage}}}
\fancyfoot[C]{\thepage}

\begin{document}
\subsubsection*{1.1 Classwork: Frequency Tables}

\begin{enumerate}[itemsep=0.15cm]

%--- Problem 1: Tally and frequency table from categorical raw data [6 marks] ---
%--- Source: Original (freshly authored, AA SL 4.1); style after Haese Core Topics SL, Ex 11D ---
\item[1.]

Twenty-four students were asked how they usually travel to school.
The replies were recorded in the order they were given:

{\setlength{\tabcolsep}{4pt}
\begin{center}
\begin{tabular}{cccccccc}
bus & subway & walk & car & subway & walk & bike & subway \\
walk & bus & subway & subway & car & walk & subway & bus \\
subway & bike & walk & bus & car & subway & walk & bus \\
\end{tabular}
\end{center}}

\begin{enumerate}[itemsep=0.05cm]
  \item Complete the tally and frequency table below.
  \item How many students were surveyed altogether?
  \item Write down the mode of the data.
  \item What fraction of the students walk to school? Give your answer
  as a fraction or as a percentage.
  \item Is \emph{how you travel to school} a categorical variable or a
  quantitative variable? Answer in one sentence.
\end{enumerate}

{\renewcommand{\arraystretch}{2.0}
\begin{center}
\begin{tabular}{|l|p{6cm}|c|}
\hline
\textbf{How you travel} & \hfil \textbf{Tally} & \textbf{Frequency} \\ \hline
walk   & & \\ \hline
bus    & & \\ \hline
subway & & \\ \hline
car    & & \\ \hline
bike   & & \\ \hline
\end{tabular}
\end{center}}

\begin{tikzpicture}
  \draw (-0.6,0) rectangle (14.6,6.2);
  \foreach \y in {1,2,...,5} \draw[dotted] (0,\y)--(14.1,\y);
\end{tikzpicture}

\newpage

%--- Problem 2: Frequency table from discrete numerical raw data [6 marks] ---
%--- Source: Original (freshly authored, AA SL 4.1); style after Haese Core Topics SL, Ex 11D ---
\item[2.]

Each of the 25 students in a class was asked how many brothers and
sisters they have. The replies were:

\begin{center}
\begin{tabular}{ccccc}
1 & 2 & 0 & 3 & 1 \\
1 & 2 & 4 & 1 & 2 \\
1 & 3 & 0 & 2 & 1 \\
1 & 2 & 3 & 1 & 2 \\
1 & 4 & 0 & 2 & 3 \\
\end{tabular}
\end{center}

\begin{enumerate}[itemsep=0.05cm]
  \item Complete the frequency table below.
  \item How many students are in the class?
  \item Write down the mode of the data.
  \item How many students have at least 2 brothers and sisters?
  \item Is \emph{number of brothers and sisters} discrete or continuous?
  Give a short reason.
\end{enumerate}

{\renewcommand{\arraystretch}{2.0}
\begin{center}
\begin{tabular}{|c|p{6cm}|c|}
\hline
\textbf{Number of siblings} & \hfil \textbf{Tally} & \textbf{Frequency} \\ \hline
0 & & \\ \hline
1 & & \\ \hline
2 & & \\ \hline
3 & & \\ \hline
4 & & \\ \hline
\end{tabular}
\end{center}}

\begin{tikzpicture}
  \draw (-0.6,0) rectangle (14.6,6.9);
  \foreach \y in {1,2,...,6} \draw[dotted] (0,\y)--(14.1,\y);
\end{tikzpicture}

\end{enumerate}
\end{document}
